\documentclass[openany]{ctexbook}
\usepackage{amsmath,amssymb}
\usepackage[hidelinks]{hyperref}
\newcommand{\note}[1]{\footnote{译注：#1}}
\title{代数几何基础第五卷预备笔记\\9.4编排提示、9.5残文及10.1}
\author{据用户上传英译编辑稿译出}
\date{2026年11月15日}
\begin{document}
\maketitle
\chapter*{范围与校订状态}
本批连续翻译底本第38页结合8.18的引导句起，至第39页10.1最后一段止，保留原脚注54、55。9.4在底本中尚待陈述，9.5只有残文；本批不补造其结论。10.1中关于例外集与其闭包的措辞存在编辑疑点，单独标明，不把草稿断言作为已独立证明的结果。下一处为第39页10.2“By a polarity”。本批为独立增补，未合入原合订本；此前8.8缺文及9.3的疑点均未消除。
\chapter*{9.4编排提示与9.5残文}
结合8.18，我们得到一个推论9.4。（〔英译者〕待陈述。）

最后，还须将后者与9.2结合，以便在“优良情形”下得到一个汇总性的定理。

\noindent\textbf{定理9.5（第52页或第55页字迹不清）.}\footnote{原脚注54：“请询问A.G.，我未能理解。〔英译者〕”}

〔残文〕（若 $X$ 在 $k$ 上光滑、固有且几何整，且 $f:X\to P$ 非分歧，$X$ 的维数 $\ge2$（请询问Grothendieck），通过考虑结果 ???，当 $x\to p$ 是一个浸入时（$Y_\xi$ 的奇异点簇。\note{底本这段语法与括号均未闭合，实际印为小写 $x\to p$；此处保留而不据猜测补写定理。所谓第52或55页是底本自身的编排标记，不是本批PDF页码。}
\chapter*{§10 超平面截面束与簇的爆破的纤维化}
\noindent\textbf{10.1.}
设 $Z$ 是 $P^v$ 中相对于一个可构造性质 $\mathcal P$ 的 $\mathcal P$-例外集，使得 $Z$ 是 $P^v$ 的可构造子集。设 $S=\operatorname{Spec}(k)$。我们将看到（参见第12节，那里将补述一些无疑本应在前几节出现的内容），要有
\[
\operatorname{codim}(Z,P^v)\ge2,
\]
〔以下为原稿断言，其中单条直线的充分条件待核。〕必要且充分的是：$P^v$ 中“每一条充分一般的直线” $L$ 都不与 $Z$ 相交，或者也可以说（甚至）不与 $Z$ 相交；而只要存在一条（单独一条）不与 $Z$ 相交的 $L$ 就足够了〔底本编辑批注：“应为 $Z$；A.G.的错误，P.B.”〕。\footnote{原脚注55：“请询问A.G.”}\note{这一整句有未解决的排印或编辑问题：底本在“or even”前后及“should be”前后均印作 $Z$，未见可辨的上划线，而下一句确实印为 $\bar z$。这里按实际可辨文字保留重复，不把它擅自恢复成关于 $\overline Z$ 的定理。对一般可构造集，“存在一条直线不交 $Z$”不足以推出余维至少2，例如在射影平面中从一条直线删去一点所得的集合；经过该点的另一条直线可避开该集合。故该处不能被当作已核实的充分条件。底本此处及后文数次写 $L\subset P$，但开头及10.2明定参数直线在对偶空间 $P^v$ 中；正文在这些明确的参数空间处统一为 $P^v$。}
若 $k$ 无限，则必要且充分的是存在一条不与 $\overline Z$ 相交的直线 $L\subset P^v$。\note{底本印作小写 $\bar z$；上下文没有引入另一个 $z$，此处按同一例外集统一为 $\overline Z$。这是原稿对上述余维条件的表述，未在本批中独立证明；“充分一般”及直线的定义域约定仍须结合第12节核查。}
我们把 $L$-预概形 $Y_L$ 称为由 $P^v$ 中直线 $L$ 定义的 $X$ 的一个超平面截面线性束（此定义对任意 $S$ 都有效）。

于是，将上述考虑与第8、9节的结果结合，首先在 $S$ 为一个无限基域的谱的情形下，便得到这样的束的存在性判据：所有纤维 $Y_\xi$（$\xi\in L$）都满足性质 $\mathcal P$。考虑到8.2，若对 $X$ 的每一个伴随素循环 $T$ 都有
\[
\dim f(T)>0,
\]
则可以要求束 $Y_L$ 在 $L$ 上平坦，办法是取性质
\[
\mathcal P'=\mathcal P+\text{$\phi_\xi$ 的正则性条件}.
\]
在 $S$ 任意的情形下，若 $k(s)$ 无限，而且已知 $Z$ 为闭集，则仍可按照7.1的方法，在 $S$ 的给定点 $s$ 的某个开邻域上构造这样的束。（第5节的结果以及 $X\to S$ 固有的假设，在各种不同情形下保证上述闭性。）

为了妥善处理此事，宜在作出这种一般性说明后，给出若干汇总性陈述，将前面的结果实际应用于若干这类性质（也包括模的性质）。至少，必须在这里用线性束的语言重新表述9.5；这是一项在几何应用中经常使用的事实。\note{此句是原稿的写作计划及评语，并不意味着底本中缺失的9.5已经补齐。}
\end{document}
